\documentclass[11pt, a4paper]{article}
\usepackage[margin=2.5cm]{geometry}
\usepackage{amsmath, amssymb}
\usepackage{mathtools}
\usepackage{bm}
\usepackage{hyperref}
\usepackage{xcolor}
\usepackage{booktabs}
\usepackage{enumerate}

\hypersetup{colorlinks=true, linkcolor=blue, urlcolor=blue}

\title{\textbf{IFT-Guided Profile Likelihood Scan}\\[0.5em]
\large Algorithm Description}
\author{Newtrinos.jl}
\date{}

\newcommand{\nupar}{\bm{\nu}}
\newcommand{\zz}{\mathbf{z}}
\newcommand{\Sig}{\boldsymbol{\Sigma}}
\newcommand{\Lam}{\boldsymbol{\Lambda}}

% Simple framed pseudocode using standard tabbing
\newenvironment{pseudocode}[1]{%
  \bigskip
  \noindent\textbf{Algorithm:} #1\\[-0.3em]
  \noindent\rule{\linewidth}{0.8pt}\\[-0.5em]
  \noindent\begin{enumerate}
}{%
  \end{enumerate}
  \noindent\rule{\linewidth}{0.4pt}
  \bigskip
}

\begin{document}
\maketitle

\section*{Motivation}

Profile likelihood scans require solving a separate optimisation at each
grid point $\vartheta_k$:
\begin{equation}
  \hat{\nupar}(\vartheta_k)
  \;=\; \arg\max_{\nupar}\;\bigl[\log\mathcal{L}(\vartheta_k, \nupar)
                                  + \log\pi(\nupar)\bigr],
\end{equation}
where $\mathcal{L}$ is the likelihood and $\pi(\nupar)$ the nuisance prior.
For experiments with $\mathcal{O}(40\text{--}90)$ nuisance parameters, each
optimisation is expensive and a poor starting point can trap the optimiser in a
local minimum.

The IFT-guided scan uses a local Gaussian (MGVI) approximation of the posterior
at the current optimum to \emph{predict} the nuisance optimum at the next grid
point, avoiding cold restarts.

\section{Prior-Normalised Coordinates ($z$-space)}

Every nuisance parameter $\nu_i$ with prior $\pi_i$ is mapped to
\begin{equation}
  z_i \;=\; \Phi^{-1}\!\bigl(F_{\pi_i}(\nu_i)\bigr),
\end{equation}
where $F_{\pi_i}$ is the prior CDF and $\Phi^{-1}$ the standard-Normal quantile
function.  Collecting all nuisance parameters gives
$\zz \in \mathbb{R}^n$ with prior $\zz \sim \mathcal{N}(\mathbf{0}, \mathbf{I})$.

The scan parameter $\vartheta_k$ is fixed to a point mass at each grid value and
is \emph{excluded} from $\zz$; it never enters the $z$-space optimisation.

At grid point $\vartheta_k$ the $z$-space posterior is
\begin{equation}
  \log p_k(\zz)
  \;=\; \log\mathcal{L}\!\bigl(\vartheta_k,\, f(\zz)\bigr)
        - \tfrac{1}{2}\|\zz\|^2 + \text{const},
\end{equation}
where $f:\mathbb{R}^n\!\to\!\mathbb{R}^n$ maps back to physical nuisance space.

\section{Local MGVI Approximation}

After optimising at grid point $\vartheta_{k-1}$ we obtain $\zz^\star_{k-1}$.
Assuming the likelihood arises from a forward model
$\bm{\mu}(\vartheta,\nupar)$ with observation-noise Fisher information
$\mathbf{F}\in\mathbb{R}^{m\times m}$, the MGVI posterior precision
at the current optimum is
\begin{equation}
  \Lam \;=\; \mathbf{J}^\top \mathbf{F}\, \mathbf{J} + \mathbf{I}_n,
  \qquad
  \mathbf{J} \;=\; \frac{\partial \bm{\mu}}{\partial \zz}
               \;\in\;\mathbb{R}^{m\times n},
\end{equation}
and the local posterior covariance is $\Sig = \Lam^{-1}$.  The Gaussian
approximation centred at the current optimum is
\begin{equation}
  q_{k-1}(\zz) \;\approx\; \mathcal{N}\!\bigl(\zz^\star_{k-1},\;\Sig\bigr).
\end{equation}

\section{IFT Predictor Step}

To predict the optimum at the \emph{next} scan point $\vartheta_k$, evaluate the
gradient of the new posterior at the old optimum:
\begin{equation}
  \mathbf{g}_k
  \;=\; \nabla_{\zz}\,\log p_k(\zz)\Big|_{\zz = \zz^\star_{k-1}}.
\end{equation}
This gradient is nonzero because shifting $\vartheta$ changes the likelihood, so
the old optimum is no longer stationary.  A single Newton step with the MGVI
Hessian approximation $(\Lam \approx \Sig^{-1})$ gives the predicted new optimum:
\begin{equation}
  \boxed{
    \zz^{\rm pred}_k
    \;=\; \zz^\star_{k-1}
          \;+\; \Sig\,\mathbf{g}_k.
  }
\end{equation}

\paragraph{Why this works.}
For a quadratic log-posterior (Gaussian likelihood, linear forward model), the
prediction is \emph{exact}: one Newton step reaches the true new optimum
regardless of the step size $\vartheta_k - \vartheta_{k-1}$, as verified by the
toy Gaussian validation script.  For non-quadratic likelihoods the prediction is
approximate but provides a far better starting point than the neighbour's
$\zz^\star_{k-1}$, especially when the nuisance optimum varies smoothly across
the scan.

\section{Polish Step}

The predicted $\zz^{\rm pred}_k$ initialises a gradient-based local optimisation
(L-BFGS) in $z$-space:
\begin{equation}
  \zz^\star_k
  \;=\; \arg\max_{\zz}\;\log p_k(\zz),
  \quad\text{initialised at }\zz^{\rm pred}_k.
\end{equation}
Starting near the true optimum, the optimiser requires far fewer gradient
evaluations than a cold start, reducing total cost and the probability of
converging to a spurious local maximum.

\section{BFS Traversal Order}

Grid points are visited in \emph{breadth-first order} starting from the grid
index nearest to a user-supplied seed (typically the global best-fit point).
This guarantees every point is warm-started from a direct neighbour whose
posterior is maximally similar, rather than from a point several steps away.
For a 2-D grid (e.g.\ $\vartheta_{1}$ vs.\ $\vartheta_{2}$), the BFS expands
outward ring by ring from the seed, visiting each point once.

\section{Algorithm Summary}

\begin{pseudocode}{IFT-guided profile likelihood scan}
  \item Build prior-normalising transform $f$ and its inverse $f^{-1}$.
  \item Sort grid points $\{\vartheta_k\}$ in BFS order from the seed index.
  \item For each scan point $\vartheta_k$ in BFS order:
    \begin{enumerate}[(a)]
      \item Fix $\vartheta = \vartheta_k$;
            define $\log p_k(\zz)= \log\mathcal{L}(\vartheta_k, f(\zz))
            - \tfrac12\|\zz\|^2$.
      \item \textit{If} $k$ is the BFS root: set
            $\zz_{\rm init} = f^{-1}(\bm{\theta}_{\rm seed})$.
      \item \textit{Else}:
        \begin{enumerate}[(i)]
          \item Retrieve parent optimum $\zz^\star_{\rm parent}$.
          \item Compute MGVI covariance $\Sig = \bigl(\mathbf{J}^\top
                \mathbf{F}\,\mathbf{J}+\mathbf{I}\bigr)^{-1}$
                at $\zz^\star_{\rm parent}$ under $p_{\rm parent}$.
          \item Evaluate IFT gradient
                $\mathbf{g}_k = \nabla_\zz \log p_k(\zz^\star_{\rm parent})$.
          \item Set $\zz_{\rm init} = \zz^\star_{\rm parent}
                + \Sig\,\mathbf{g}_k$.
        \end{enumerate}
      \item Polish: $\zz^\star_k =
            \text{L-BFGS}\!\bigl(\log p_k,\;\zz_{\rm init}\bigr)$.
      \item Store $\hat{\nupar}(\vartheta_k) = f(\zz^\star_k)$ and
            $\hat\ell(\vartheta_k) = \log\mathcal{L}(\vartheta_k,
            \hat{\nupar}(\vartheta_k))$.
    \end{enumerate}
\end{pseudocode}

\section{Computational Cost}

Let $n$ be the number of nuisance parameters and $m$ the data dimension.

\begin{itemize}
  \item \textbf{MGVI covariance} ($\mathcal{O}(mn^2 + n^3)$ per grid point):
    one forward-model Jacobian and one matrix inversion.
  \item \textbf{IFT gradient}: one additional likelihood gradient evaluation
    ($\mathcal{O}(mn)$ with reverse-mode AD or $\mathcal{O}(mn/B)$ batches
    with ForwardDiff chunk size $B$).
  \item \textbf{Polish}: typically $5$--$30\times$ fewer L-BFGS iterations than
    a cold start, because $\zz^{\rm pred}_k$ is already near the optimum.
\end{itemize}

For $n \gtrsim 20$ (where ForwardDiff requires multiple chunked passes) the
reduction in polish iterations typically outweighs the MGVI overhead, giving a
net speed-up of $2$--$5\times$ over sequential warm-starting alone.

\end{document}
